\chapter{Implementation}
\label{chap:implementation}

\section{Software Structure}

The prototype is implemented in Python with PyTorch, Hugging Face Transformers, Datasets,
Evaluate, and PEFT. The repository is divided into the modules in \Cref{tab:software-files}.

\begin{table}[H]
\centering
\caption{Main implementation files.}
\label{tab:software-files}
\begin{tabularx}{\textwidth}{@{}lX@{}}
\toprule
File & Responsibility \\
\midrule
\texttt{train.py} & Dataset loading, model construction, PEFT configuration, training, and evaluation.\\
\texttt{dataset\_analyzer.py} & Label, token, and embedding statistics.\\
\texttt{rank\_policy.py} & Size, distribution, complexity, and bounded rank functions.\\
\texttt{stability\_callback.py} & Importance logging and calls to PEFT's AdaLoRA allocator.\\
\texttt{run\_experiments.py} & Multi-method, multi-budget, multi-seed execution and aggregation.\\
\bottomrule
\end{tabularx}
\end{table}

\section{Dataset Loading and Preprocessing}

The implementation loads a Hugging Face dataset and uses its training split for low-resource
subsampling. If the requested budget is smaller than the full split, the training data is shuffled
using the experiment seed and the first $N$ examples are selected. Evaluation uses the complete
validation split when available. Different seeds therefore affect both model initialization and
the selected low-resource subset.

Inputs are tokenized with truncation at 256 tokens. Dynamic padding is applied per minibatch.
For classification, the number of labels is inferred from the sampled training labels. This
implementation path supports classification but does not yet handle GLUE regression tasks such as
STS-B.

\section{Model Construction}

The default model is \qwenmodel, loaded as a sequence classifier. Qwen2.5 is a family of
decoder-only language models released at multiple scales \citep{yang2024qwen25}. The 0.5B model
keeps the undergraduate experiment computationally manageable while retaining modern Transformer
attention and gated MLP projections.

For standard LoRA and AdaLoRA, Qwen targets \texttt{q\_proj} and \texttt{v\_proj}. Star-LoRA
computes the profile before attaching PEFT adapters. DistilBERT-style \texttt{q\_lin} and
\texttt{v\_lin} names are recognized, although all reported experiments use Qwen.

\section{Training Schedule}

The estimated number of optimization steps is
\begin{equation}
  T = \left\lceil\frac{N}{B}\right\rceil E,
\end{equation}
where $B$ is batch size and $E$ is the number of epochs. For standard AdaLoRA, the implementation
uses target rank 8, initial rank 12, initial and final schedule fractions of 0.1 and 0.5, and an
allocation interval near $0.02T$. Star-LoRA uses initial and final fractions of 0.2 and 0.85 and
an interval near $0.05T$.

The allocation callback invokes PEFT's \texttt{update\_and\_allocate} after optimizer steps.
The stability callback logs parameter scores before optimizer updates. The callbacks are separate:
logged stable scores do not alter the PEFT allocation decision.

\section{Recorded Artifacts}

Each run records the method, rank settings, target modules, dataset statistics for dataset-aware
runs, evaluation metrics, training metrics, random seed, model name, dataset name, and estimated
steps. Per-step raw and stable importance values are written to a CSV file.

One reproducibility limitation is that archived \texttt{experiment\_info.json} files do not store
every command-line hyperparameter, notably learning rate. Folder names and experiment commands
provide context, but future runs should serialize the complete argument dictionary, software
versions, hardware identity, and git commit.

\section{LLM-Assisted Prototype}

An experimental branch supplied dataset statistics and model structure to an external LLM and
requested a layer-wise target-module map. The Claude artifacts record
\texttt{layerwise\_policy\_source = claude} and preserve the returned map. The selected global
target rank remains 9 with initial rank 18 in the reported 5,000-example runs; the LLM primarily
changes module placement rather than replacing the numeric rank formula.

The Gemini-labeled pilot directories require special handling. Their metadata records an HTTP 429
quota failure, a null Gemini response, and \texttt{layerwise\_policy\_source = heuristic}. They
therefore duplicate the heuristic policy and cannot be treated as evidence about Gemini's
allocation quality.

\section{Reproducibility Status}

\begin{table}[H]
\centering
\caption{Reproducibility status of the current prototype.}
\label{tab:repro-checklist}
\begin{tabularx}{\textwidth}{@{}lX@{}}
\toprule
Item & Status \\
\midrule
Model and dataset & Stored in each experiment metadata file.\\
Sample budget and seed & Stored and included in output directory names.\\
Rank and target modules & Stored for each PEFT run.\\
Metrics and runtime & Stored and aggregated into CSV summaries.\\
Full CLI arguments & Partial; learning rate is not serialized in old artifacts.\\
Package versions & Not stored in the current artifacts.\\
Hardware identity & Not stored in the current artifacts.\\
LLM response provenance & Stored for Claude; failed Gemini requests and fallback are stored.\\
\bottomrule
\end{tabularx}
\end{table}
